\begin{abstract}
Label-free self-play trains a model on problems it is unsure about: problems it wrote itself on which its sampled answers disagree. R-Zero credits this rule with filtering out ill-posed problems. We show that it keeps what training has not yet resolved: among model-written problems, mostly the ill-posed ones. With Qwen3-4B/8B, Phi-4-mini and a Qwen3-14B round, the ill-posed share of admitted problems rises by 11.5--27 points within two rounds, also under blind re-judging, an open-weight judge and expert human annotation. The official R-Zero code drifts too, mostly because its challenger, rewarded with the same uncertainty signal, writes more ill-posed questions. Scoring one fixed pool with successive solvers isolates the filter: solvable problems sharpen out through the band's upper edge while problems whose text supports several readings stay. Sharpening also changes how the solver answers such problems: in Qwen3-4B, self-consistent answers to unanswerable questions rise from 28\% to 41--48\%, and stated assumptions and verbal declines both become more common, while benchmark accuracy plateaus. Most of this rise is sharpening itself; at matched curriculum size, the ill-posed problems add none that we can detect. No label-free admission rule we tested stops the drift, but offering abstention at test time works. We pre-registered 136 predictions, report every outcome, and mark exploratory analyses.
\end{abstract}
